\begin{itemize}
  \item ``\gls{sem} is \dots ideally suited for second-order problems and it is highly accurate when applied to fairly regular problems.''
  \item Tensor products are especially useful in this method for higher dimensions.
  \item $A_N=SP_{2N-1}$
  \item The use of \An is that it gives a set of second-order PDEs, which is used nicely with the \gls{sem}.
  \item Large $N$ is a problem because of strong coupling, like \Pn.
  \item The \An expansion gives a form that is similar to the even- and odd- parity and \gls{saaf} equations with a mean free path term in the now-elliptic streaming operation. $-\divergence \left[\frac{1}{\Sigma_t}\grad\left(\mathcal{A}\mathbf{f}\right)\right]$
  \item ``The main limitation deriving from the use of the continuous Galerkin approach is the fact that all elements are required to have the same number of degrees of freedom.''
  \item The Natelson benchmark was used, which seems like a neat 2D benchmark that has all reflective boundaries, so the weak form naturally satisfies these conditions (no boundary condition shenanigans). That said, no ``true'' solutions exist in the literature, though it looks like the problem has been used multiple times\dots
  \item Next benchmark was the IAEA EIR-2 benchmark, which looks like a fantastic benchmark! This one would likely show ray effects for \Sn.
  \item \gls{sem} $>$ \gls{fem} in comparison calculations (by a long-shot). Highest numerical accuracy in the shortest amount of time (1 minute vs. 1 hour).
\end{itemize}